\scanpage{53}
Nous préférons ne pas nous donner de métrique
d'avance, et prendre les autom. affines de $E$
qui conservent $E$, et exiger qu'ils \unsure{sont|soit|soient} transitifs
sur l'ens des repères. On \unsure{récupèrera|récupérera|récupère} la métrique
après-coup -- de façon unique modulo homothétie
-- mais parfois elle sera dégénérée ! (Si on exigeait
dès le début une métrique non dégénérée, on perdrait
des cas \deleted{\unsure{d'un grand intérêt}} \inserted{fort intéressants} du point de
vue arithmétique -- en fait, les seuls cas vraiment
nouveaux !) Pratiquement, ce nouveau pt
de vue consiste à regarder \uline{tt} \uline{triangle} comme
polygône régulier, ainsi que \uline{tout} \uline{parallélogramme} !
(Mais on \unsure{réc.|vis.|récup.|vrais} régulier \uline{pour} une métrique
\uline{donnée}\ldots)

\uline{Définition en forme}

1) Polygône régulier (\uline{NB} pour tt polygône, l'ens des
sommets engendre affinement $E$, donc un autom
affin est connu quand on connaît sa restriction à $S$)

2) Réalisation géométrique régulière d'un polygône
combinatoire $\Pi$\ldots
% archive sheet number: 52
